\section{Memory-Aware Chunking}
\label{sec:methodology}

The memory-aware chunking problem constitutes an optimization problem.
Given a three-dimensional volume $V=[d_1,d_2,d_3]$ representing the input data dimensions, a worker memory budget $M$ (typically 80\% of physical \ac{RAM} to allow for system overhead), and an operator \textit{op} from the computational workflow, the objective is to derive the largest cubic chunk dimension $c$ that guarantees execution without \ac{OOM} failures.

The cubic constraint simplifies the search space while maintaining good cache locality, since non-cubic chunks often lead to strided memory access patterns that degrade performance.
The optimization balances multiple objectives: maximizing chunk size to reduce scheduling overhead, ensuring memory safety with appropriate margins, and maintaining divisibility constraints so chunks evenly partition the volume.

\subsection{Memory Model}

Profiling reveals near-linear scaling between input volume and peak memory for tensor operators, yielding the model
\begin{equation}
M(V)=\alpha V+\beta f(V)+\gamma,\label{eq:model}
\end{equation}
where $V=d_1d_2d_3$, $f(V)$ captures operator-specific expansions, and $\alpha,\beta,\gamma$ represent learned parameters.
For the operators studied, $f(V)=V$, so a first-order regression suffices.
Experiments demonstrate that training on as few as 30 samples yields $R^2>0.999$, as Table~\ref{tab:feature-importance} indicates.

\begin{table}[htb]
    \centering
    \caption{Feature importance and model performance for memory prediction}
    \label{tab:feature-importance}
    \begin{tabular}{lccc}
        \toprule
        \textbf{Operator} & \textbf{$R^2$} & \textbf{RMSE (MB)} & \textbf{Training Samples} \\
        \midrule
        Envelope & 0.9995 & 0.82 & 30 \\
        Gaussian Filter & 0.9997 & 0.64 & 30 \\
        \ac{GST3D} & 0.9993 & 1.23 & 30 \\
        \bottomrule
    \end{tabular}
\end{table}

\subsection{Chunk-Size Algorithm}

Algorithm~\ref{alg:memaware} computes $c$ in four steps: 
(1) predict peak memory for the full volume; 
(2) derive per-voxel cost; 
(3) compute the largest cube that respects a safety factor $s$ (default $0.8$); and 
(4) decrement $c$ until it divides all dimensions.
The procedure runs in $\mathcal{O}(d_{\max})$ time, negligible relative to operator execution.

\begin{algorithm}[tb]
\caption{Memory-aware chunk sizing}
\label{alg:memaware}
\begin{algorithmic}[1]
\REQUIRE Volume $V=[d_1,d_2,d_3]$, memory limit $M$, safety factor $s$, predictor $\mathcal{M}$
\ENSURE Chunk dimension $c$
\STATE $\textit{peak}\leftarrow\mathcal{M}.\text{predict}(V)$
\STATE $\textit{cost}\leftarrow\textit{peak}/(d_1d_2d_3)$
\STATE $c\leftarrow\lfloor ((M\,s)/\textit{cost})^{1/3}\rfloor$
\WHILE{$\exists i\, :\, d_i\bmod c\neq0$ \textbf{and} $c>1$}
    \STATE $c\leftarrow c-1$
\ENDWHILE
\RETURN $c$
\end{algorithmic}
\end{algorithm}

\subsection{Integration with \dask}

The algorithm integrates seamlessly with \dask's lazy evaluation model through its chunk suggestion \ac{API}.
When users create a \dask array or apply an operator, the system intercepts the chunking decision point.
At graph-construction time, before any computation begins, the function inspects the input array shape, identifies the operator being applied, and queries the corresponding predictor.

The implementation leverages \dask's plugin architecture as follows:
\begin{lstlisting}[language=Python]
def memory_aware_chunks(shape, dtype, operator_name):
    predictor = load_predictor(operator_name)
    memory_limit = get_worker_memory()
    return compute_chunk_size(shape, predictor, memory_limit)
\end{lstlisting}

This design requires no modifications to \dask's core scheduler or executor.
Existing code continues to work unchanged; setting a configuration flag activates memory-aware chunking.
The system stores predictors as lightweight pickle files (typically \textless 1~KB each) that load in milliseconds, adding negligible overhead to graph construction.

\subsection{Benefits}

The memory-aware chunking method provides several key benefits that address the limitations of current approaches:

\textbf{(i) Memory Safety Guarantees:} By modeling peak memory usage including all intermediate allocations, the system prevents OOM failures that plague production deployments.
The safety factor provides a tunable margin for system overhead and measurement uncertainty.

\textbf{(ii) Elimination of Manual Tuning:} The system eliminates the need for researchers to iteratively search for viable chunk sizes, a process that can waste days of compute time on large clusters.
The system automatically adapts to different hardware configurations and operator characteristics.

\textbf{(iii) Performance Optimization:} The algorithm selects the largest safe chunk size, thereby minimizing scheduling overhead and maximizing cache efficiency when memory permits.
This approach differs from conservative methods that use unnecessarily small chunks "to be safe."

\textbf{(iv) Operator Awareness:} Unlike fixed-size heuristics, the method adapts to each operator's specific memory profile.
Memory-intensive operations like \ac{GST3D} receive smaller chunks than lightweight operations like element-wise multiplication.

\textbf{(v) Practical Deployment:} The approach requires minimal training data (30-50 samples), integrates transparently with existing code, and adds negligible runtime overhead (\textless 0.1\% of typical operator execution time).